%% =====================================================================
%%  BEAMER POSTER TEMPLATE -- A0 portrait, 3 columns, academic style
%% =====================================================================
%%  HOW TO USE
%%   1. Compile with  pdflatex  (run twice if you use \ref / \cite).
%%   2. Search for  "EDIT:"      -> lines you will normally change.
%%      Search for  "OPTIONAL:"  -> things you may delete or switch off.
%%   3. Put hbni_logo.png, logo2.png and your figures in the SAME folder
%%      as this file.  A missing image is replaced by a grey placeholder
%%      box, so the file always compiles.
%%
%%  LAYOUT
%%      [ logo 1 |  TITLE  /  AUTHORS  /  AFFILIATIONS  | logo 2 ]
%%      [ column 1 |  column 2  |  column 3 ]
%%      [ take-home bar ]
%%      [ references | acknowledgements | QR code ]
%%      [ contact | conference ]
%%
%%  IF TEXT OVERFLOWS (a quick guide)
%%   - Title too long / wraps to 4 lines -> lower \TitleSize (header section).
%%   - Column content too long           -> lower "scale" below (1.5 -> 1.3),
%%                                          or move a block to another column.
%%   - Wide table sticks out of a block  -> wrap it in \resizebox{\linewidth}{!}{...}
%% =====================================================================

%% Beamer already loads xcolor, so the 'table' option (for \rowcolor) must be
%% passed BEFORE \documentclass. This avoids the "Option clash for xcolor" error.
\PassOptionsToPackage{table,dvipsnames}{xcolor}
\documentclass[final]{beamer}

%% ---------------------------------------------------------------------
%%  PAGE SIZE  (EDIT)
%%  size = a0 | a1 | a2 | a3      orientation = portrait | landscape
%%  scale = global text magnification (bigger = larger text).
%% ---------------------------------------------------------------------
\usepackage[orientation=portrait,size=a0,scale=1.5]{beamerposter}
%  Change Scale for font size in the main body

%% ---------------------------------------------------------------------
%%  PACKAGES  (all ship with TeX Live / MiKTeX / Overleaf)
%% ---------------------------------------------------------------------
\usepackage[utf8]{inputenc}
\usepackage[T1]{fontenc}
%% ---- FONT CHOICE (EDIT): keep exactly ONE of the two options below ----
%% Option A (default): serif, Palatino text + matching maths
\usepackage{mathpazo}
%% Option B: sans-serif, Helvetica text. To use it, put % in front of the
%% mathpazo line above, remove the % from the 3 lines below, and also
%% change \usefonttheme{serif} further down to \usefonttheme{professionalfonts}.
% \usepackage[scaled=0.95]{helvet}
% \renewcommand{\familydefault}{\sfdefault}
% \usepackage{sansmath}\sansmath   % sans-serif maths to match
\usepackage{graphicx}
\usepackage{booktabs}          % professional tables
\usepackage{amsmath,amssymb}
\usepackage{xcolor}
\usepackage[most]{tcolorbox}   % coloured call-out boxes, header, footer
\usepackage{qrcode}            % OPTIONAL: QR code (delete with its use in the footer)
\usepackage{microtype}         % better spacing, fewer overfull lines

%% ---------------------------------------------------------------------
%%  COLOURS: "Oxford Navy + Burgundy"  (EDIT the RGB numbers to re-theme)
%%  Rule of thumb: ONE dark colour, ONE accent, everything else neutral.
%% ---------------------------------------------------------------------
\definecolor{Primary}  {RGB}{ 11, 37, 69}   % Oxford navy    -> header, block titles
\definecolor{Secondary}{RGB}{ 62, 92,118}   % slate blue     -> secondary blocks
\definecolor{Accent}   {RGB}{164, 36, 59}   % burgundy       -> emphasis, warnings, rules
\definecolor{Panel}    {RGB}{241,244,248}   % pale grey-blue -> block body / info boxes
\definecolor{Rule}     {RGB}{201,211,223}   % light grey-blue-> thin borders
\definecolor{TextCol}  {RGB}{ 26, 26, 26}   % near black      -> body text

%% ---------------------------------------------------------------------
%%  BEAMER LOOK  (rarely needs editing)
%% ---------------------------------------------------------------------
\usetheme{default}
\usefonttheme{serif}                                     % serif text in all beamer parts
\setbeamertemplate{navigation symbols}{}
\setbeamercolor{background canvas}{bg=white}
\setbeamercolor{normal text}{fg=TextCol}
\usebeamercolor[fg]{normal text}

\setbeamercolor{block title}         {fg=white,bg=Primary}
\setbeamercolor{block body}          {fg=TextCol,bg=Panel}
\setbeamercolor{block title alerted} {fg=white,bg=Accent}
\setbeamercolor{block body alerted}  {fg=TextCol,bg=Panel}
\setbeamercolor{block title example} {fg=white,bg=Secondary}
\setbeamercolor{block body example}  {fg=TextCol,bg=Panel}

\setbeamerfont{block title}{size=\large,series=\bfseries}
\setbeamerfont{block body}{size=\normalsize}

% Square-cornered block: navy title bar + pale body. Shared by all 3 block types.
\newcommand{\PosterBlockBegin}[2]{%
  \par
  \begin{beamercolorbox}[wd=\linewidth,leftskip=0.7em,sep=0.8ex]{#1}%
    \strut\usebeamerfont*{block title}\insertblocktitle
  \end{beamercolorbox}%
  \par\nointerlineskip
  \usebeamerfont{block body}%
  \begin{beamercolorbox}[wd=\linewidth,sep=0.7em,vmode]{#2}%
    \setlength{\linewidth}{\dimexpr\linewidth-1.4em\relax}\setlength{\columnwidth}{\linewidth}\setlength{\textwidth}{\linewidth}% figures/boxes = inner width, never overflow
    }
\newcommand{\PosterBlockEnd}{\par\end{beamercolorbox}\vskip 8pt}
\setbeamertemplate{block begin}          {\PosterBlockBegin{block title}{block body}}
\setbeamertemplate{block alerted begin}  {\PosterBlockBegin{block title alerted}{block body alerted}}
\setbeamertemplate{block example begin}  {\PosterBlockBegin{block title example}{block body example}}
\setbeamertemplate{block end}            {\PosterBlockEnd}
\setbeamertemplate{block alerted end}    {\PosterBlockEnd}
\setbeamertemplate{block example end}    {\PosterBlockEnd}

% Bullets and numbering
\setbeamertemplate{itemize item}{\raisebox{0.2ex}{\tiny$\blacksquare$}}
\setbeamertemplate{itemize subitem}{\raisebox{0.1ex}{\tiny$\blacktriangleright$}}
\setbeamertemplate{enumerate item}{\insertenumlabel.}
\setbeamercolor{itemize item}{fg=Accent}
\setbeamercolor{itemize subitem}{fg=Secondary}
\setbeamercolor{enumerate item}{fg=Accent}

% TIGHT SPACING (this removes wasted white space)
\addtobeamertemplate{itemize body begin}  {\itemsep=3pt \topsep=2pt \parsep=0pt }{}
\addtobeamertemplate{enumerate body begin}{\itemsep=3pt \topsep=2pt \parsep=0pt }{}
\setbeamersize{text margin left=1.5cm,text margin right=1.5cm}
\setlength{\parskip}{4pt}
\setlength{\parindent}{0pt}
\setlength{\abovedisplayskip}{4pt}   \setlength{\belowdisplayskip}{4pt}
\setlength{\abovedisplayshortskip}{2pt} \setlength{\belowdisplayshortskip}{2pt}

%% ---------------------------------------------------------------------
%%  HANDY MACROS
%% ---------------------------------------------------------------------
\newcommand{\kw}[1]{\textbf{\textcolor{Primary}{#1}}}     % keyword in navy bold
\newcommand{\num}[1]{\textbf{\textcolor{Accent}{#1}}}     % key number in burgundy bold

% \PosterLogo{file}{max width}{max height}: image scaled to FIT the box
% (never overflows), or a grey placeholder when the file is missing.
\newcommand{\PosterLogo}[3]{%
  \IfFileExists{#1}%
   {\includegraphics[width=#2,height=#3,keepaspectratio]{#1}}%
   {\fcolorbox{white}{white!25!Primary}{\parbox[c][#3][c]{#2}{\centering\color{white}\bfseries LOGO\\[2pt]{\small\expandafter\detokenize\expandafter{#1}}}}}}

% \fig{file}{width}: image, or placeholder (height 0.19 of page) if missing
\newcommand{\fig}[2]{%
  \IfFileExists{#1}{\includegraphics[width=#2]{#1}}%
  {\fcolorbox{Rule}{white}{\parbox[c][0.19\textheight][c]{#2}{\centering\bfseries Figure placeholder\\[2pt]{\small\expandafter\detokenize\expandafter{#1}}}}}}

% Call-out boxes to use inside or between blocks
\tcbset{
  keybox/.style ={enhanced,colback=white,colframe=Accent,boxrule=1.6pt,arc=0pt,
                  left=6pt,right=6pt,top=4pt,bottom=4pt,before=\vspace{2pt},after=\vspace{2pt}},
  infobox/.style={enhanced,colback=white,colframe=Rule,boxrule=1pt,arc=0pt,
                  borderline west={5pt}{0pt}{Secondary},
                  left=10pt,right=6pt,top=6pt,bottom=6pt,before=\vspace{2pt},after=\vspace{2pt}},
}

%% =====================================================================
%%  POSTER INFORMATION  --  EDIT EVERYTHING IN THIS SECTION
%% =====================================================================
\newcommand{\PosterTitle}    {Your Poster Title Goes Here:}
\newcommand{\PosterSubtitle} {A Concise Subtitle Describing the Contribution}   % OPTIONAL: leave {} to hide
\newcommand{\Authors}        {First Author$^{1,2}$ \quad Second Author$^{1,2}$ \quad Third Author$^{2}$}
\newcommand{\Affiliations}   {$^{1}$Department of X, Institute A, City, Country \quad $^{2}$Institute B, City, Country}
\newcommand{\ContactEmail}   {first.author@institute.edu \;$\cdot$\; second.author@institute.edu}
\newcommand{\EventName}      {Conference / Workshop Name 2026}

%% Logos (png / jpg / pdf).  Each is scaled to fit inside Width x Height.
\newcommand{\LeftLogo}       {media/hbni_logo.png}
\newcommand{\RightLogo}      {media/Niser_logo.png}
\newcommand{\LogoWidth}      {0.13\paperwidth}   % max width of one logo
\newcommand{\LogoHeight}     {8cm}               % max height of one logo

%% Header font sizes (pt).  Lower them if a long title does not fit.
\newcommand{\TitleSize}      {\fontsize{84}{98}\selectfont}
\newcommand{\SubtitleSize}   {\fontsize{46}{54}\selectfont}
\newcommand{\AuthorSize}     {\fontsize{48}{56}\selectfont}
\newcommand{\AffilSize}      {\fontsize{34}{40}\selectfont}

%% =====================================================================
\begin{document}
\begin{frame}[t]
\vspace{-1.4cm}

%% ---------------------------------------------------------------------
%%  HEADER: logo | title + authors | logo
%%  Widths of logo + title + logo must fit; the title box is 0.66 of the line.
%% ---------------------------------------------------------------------
\begin{tcolorbox}[enhanced,width=\linewidth,colback=Primary,colframe=Primary,
  arc=0pt,boxrule=0pt,left=1cm,right=1cm,top=0.7cm,bottom=0.7cm,
  borderline south={14pt}{0pt}{Accent}]   % OPTIONAL: burgundy rule under header
  \begin{minipage}[c]{\LogoWidth}\centering
    \PosterLogo{\LeftLogo}{\LogoWidth}{\LogoHeight}
  \end{minipage}\hfill
  \begin{minipage}[c]{0.66\linewidth}\centering
    {\color{white}\bfseries\TitleSize \PosterTitle\par}
    \vspace{0.25cm}
    {\color{white!80!Primary}\itshape\SubtitleSize \PosterSubtitle\par}   % OPTIONAL subtitle
    \vspace{0.45cm}
    {\color{white}\AuthorSize \Authors\par}
    \vspace{0.15cm}
    {\color{white!78!Primary}\AffilSize \Affiliations\par}
  \end{minipage}\hfill
  \begin{minipage}[c]{\LogoWidth}\centering
    \PosterLogo{\RightLogo}{\LogoWidth}{\LogoHeight}
  \end{minipage}
\end{tcolorbox}

\vspace{0.2cm}

%% =====================================================================
%%  BODY: THREE COLUMNS
%%  Blocks:   \begin{block}{Title}...\end{block}       navy title
%%            \begin{alertblock}{Title}...             burgundy title
%%            \begin{exampleblock}{Title}...           slate-blue title
%%  For TWO columns use \begin{column}{0.485\textwidth} twice.
%% =====================================================================
\begin{columns}[T,onlytextwidth]

%% ------------------------- COLUMN 1 ---------------------------------
\begin{column}{0.315\textwidth}

\begin{block}{Introduction and Motivation}
  Graph neural networks are deployed in \kw{sensitive domains}, where regulators
  require both explanations and privacy. Write two or three sentences that
  explain the big picture and why it matters.
  \begin{itemize}
    \item Why does this problem matter in the real world?
    \item What is missing in existing work?
    \item What is the key idea of this work?
  \end{itemize}
  \begin{tcolorbox}[keybox]
    \centering\bfseries Research question: can an adversary recover private
    structure from public outputs?
  \end{tcolorbox}
\end{block}

\begin{block}{Problem Definition}
  Given input $x \in \mathcal{X}$, we seek parameters $\theta^{\star}$ with
  \[
    \theta^{\star} = \arg\min_{\theta}\;
    \mathbb{E}_{(x,y)}\big[\,\mathcal{L}(f_\theta(x),\,y)\,\big].
  \]
  \textbf{Assumptions:}
  \begin{enumerate}
    \item The adversary sees only public outputs.
    \item Model parameters are not accessible.
  \end{enumerate}
\end{block}

\begin{alertblock}{Limitations of Existing Work}
  \begin{itemize}
    \item Existing attacks assume clean data; accuracy drops by \num{30\%} under noise.
    \item They do not scale beyond small graphs (out of memory).
  \end{itemize}
\end{alertblock}

\begin{block}{Contributions}
  \begin{itemize}
    \item \kw{Contribution 1:} a new method that does something useful.
    \item \kw{Contribution 2:} a theoretical result supporting the method.
    \item \kw{Contribution 3:} extensive experiments on five benchmarks.
  \end{itemize}
\end{block}

\end{column}\hfill

%% ------------------------- COLUMN 2 ---------------------------------
\begin{column}{0.315\textwidth}

\begin{block}{Proposed Method}
  \centering
  \fig{architecture.pdf}{\linewidth}   % EDIT: your diagram
  \par\vspace{2pt}
  \raggedright\textbf{Figure 1.} One-sentence caption describing the pipeline.
\end{block}

\begin{exampleblock}{Key Components}
  \begin{tcolorbox}[infobox]\textbf{(1) Component A.} What it does, in one line.\end{tcolorbox}
  \begin{tcolorbox}[infobox]\textbf{(2) Component B.} What it does, in one line.\end{tcolorbox}
  \begin{tcolorbox}[infobox]\textbf{(3) Component C.} What it does, in one line.\end{tcolorbox}
  \textbf{Training objective:}
  \begin{equation}
    \mathcal{L}_{\text{total}} = \mathcal{L}_{A} + \lambda\,\mathcal{L}_{B}
  \end{equation}
\end{exampleblock}

\begin{block}{Algorithm}
  \begin{enumerate}
    \item Collect and clean the data.
    \item Train the model with the joint objective.
    \item Evaluate on the benchmark datasets.
  \end{enumerate}
\end{block}

\begin{block}{Experimental Setup}
  \textbf{Datasets:} list them here. \textbf{Metrics:} AUC and AP.
  \textbf{Baselines:} name the methods you compare with.
  \textbf{Implementation:} optimiser, learning rate, hardware.
\end{block}

\end{column}\hfill

%% ------------------------- COLUMN 3 ---------------------------------
\begin{column}{0.315\textwidth}

\begin{block}{Results}
  \centering
  \renewcommand{\arraystretch}{1.2}\setlength{\tabcolsep}{6pt}
  % Wide table? wrap the tabular in \resizebox{\linewidth}{!}{ ... }
  \begin{tabular}{lccc}
    \toprule
    \textbf{Method} & \textbf{Data 1} & \textbf{Data 2} & \textbf{Data 3} \\
    \midrule
    Baseline A & 0.52 & 0.54 & 0.51 \\
    Baseline B & 0.88 & 0.93 & 0.49 \\
    \midrule
    \rowcolor{Accent!12}
    \textbf{Ours} & \textbf{0.94} & \textbf{0.98} & \underline{0.60} \\
    \bottomrule
  \end{tabular}
  \par\vspace{3pt}
  \raggedright\textbf{Table 1.} Bold: best; underline: second best.
\end{block}

\begin{block}{Analysis}
  \centering
  \fig{results_plot.pdf}{\linewidth}   % EDIT: your plot
  \par\vspace{2pt}
  \raggedright\textbf{Figure 2.} Caption stating the trend the reader should notice.
\end{block}

\begin{block}{Discussion}
  Explain \emph{why} the results look the way they do, and mention the
  main limitation of your approach in one or two sentences.
\end{block}

\begin{block}{Conclusions}
  \begin{enumerate}
    \item \textbf{Main finding:} one sentence with a number, e.g.\ \num{+10\%} AUC.
    \item \textbf{Second finding:} one sentence.
    \item \textbf{Future work:} one sentence.
  \end{enumerate}
\end{block}

\end{column}
\end{columns}

%% ---------------------------------------------------------------------
%%  OPTIONAL: full-width take-home message (delete if not needed)
%% ---------------------------------------------------------------------
\vspace{6pt}
\begin{tcolorbox}[keybox]
  \centering\bfseries Take-home message: one sentence a passer-by can remember.
\end{tcolorbox}

%% OPTIONAL: \vfill pushes the footer to the page bottom when content is short.
\vfill

%% ---------------------------------------------------------------------
%%  FOOTER: references | acknowledgements | QR code
%% ---------------------------------------------------------------------
\begin{tcolorbox}[enhanced,colback=white,colframe=Primary,boxrule=1.4pt,arc=0pt,
  left=10pt,right=10pt,top=6pt,bottom=6pt]
  \small
  \begin{minipage}[t]{0.52\linewidth}
    \textbf{\color{Primary}References}\par\vspace{2pt}
    % Few refs: type by hand. Many refs: use BibTeX instead.
    [1] A.~Author, B.~Author. \emph{Paper title.} Venue, 2024.\par
    [2] C.~Author. \emph{Another paper title.} Venue, 2023.\par
    [3] D.~Author et~al. \emph{Third paper title.} Venue, 2022.
  \end{minipage}\hfill
  \begin{minipage}[t]{0.29\linewidth}
    \textbf{\color{Primary}Acknowledgements}\par\vspace{2pt}
    This work was supported by \emph{Funding Agency}, Grant No.~XXXX.   % EDIT
  \end{minipage}\hfill
  \begin{minipage}[t]{0.11\linewidth}\centering
    \qrcode[height=3cm]{https://smlab.niser.ac.in/}\\[-1pt]   % OPTIONAL: link to paper/code
    {\scriptsize Paper \& code}
  \end{minipage}
\end{tcolorbox}

%% ---------------------------------------------------------------------
%%  BOTTOM STRIP
%% ---------------------------------------------------------------------
\begin{tcolorbox}[enhanced,colback=Primary,colframe=Primary,arc=0pt,boxrule=0pt,
  left=1cm,right=1cm,top=0.3cm,bottom=0.3cm,before=\vspace{2pt},after=,
  borderline north={8pt}{0pt}{Accent}]
  \color{white}\normalsize
  \textbf{Contact:} \ContactEmail \hfill \textbf{\EventName}
\end{tcolorbox}

\end{frame}
\end{document}